% !TEX program = xelatex

\documentclass[11pt,a4paper]{ctexart}

% ====== 字体 ======
\usepackage{fontspec}
\setmainfont{Latin Modern Roman}
\usepackage{unicode-math}
% 修改：XITS Math 拥有原生粗体数学字形，解决 \symbfit 加粗无效的问题
\setmathfont{XITS Math}

% ====== 颜色支持 ======
\usepackage{xcolor}

% ====== 页面设置 ======
\usepackage[margin=2.5cm]{geometry}
\usepackage{setspace}
\onehalfspacing
\pagestyle{plain}

% ====== 数学宏包 ======
\usepackage{mathtools}
\usepackage{amsmath}
% \usepackage{amssymb}   % 删除
% \usepackage{bm}        % 彻底删除

% ====== 三线表 + 数字对齐 ======
\usepackage{booktabs}
\usepackage{siunitx}
\usepackage{enumitem}

% ====== 交叉引用与超链接 ======
\usepackage[colorlinks=true,
            linkcolor=blue,
            citecolor=green,
            urlcolor=cyan]{hyperref}
\usepackage[noabbrev,nameinlink]{cleveref}

% ====== 自定义命令 ======
\newcommand{\dd}{\mathrm{d}}
\newcommand{\source}[1]{\par \hfill ——#1}

\title{电动力学讨论：“囚禁”电场线}
\author{王晨暄}
\date{\today}

\begin{document}
\maketitle

{
    \kaishu 是否存在一个非球对称的电荷分布体系$\rho(\symbfit{r})$，使得远场来看，
    体系的电$2^l$极矩的$2l+1$个分量全部为零，其中$l=1,2,3\cdots$。
    \label{question}
}

\bigskip

这个问题十分有趣。一方面，如果存在这样的电荷分布，那么在远场做电势函数的多级展开，则有$\phi=C$，其中$C$为常数。
也就是说：$\exists R>0,\,\forall\left|\symbfit{r}\right|>R,\,\symbfit{E}(\symbfit{r})=-\nabla\phi=0$，即：
电场线被“囚禁”在有限的区域内，而没有电场线延伸到无穷远处。另一方面，电场的多级展开可以让我们将
有限分布电荷的电场展开成级数形式。较为简单的一种电$2^n$极矩的第$(i_1,i_2,\cdots,i_n)$分量定义为：

\begin{equation}
    Q_{i_1i_2\cdots i_n}=\int\rho(\symbfit{r'})r'_{i_1}r'_{i_2}\cdots r'_{i_n}\dd^3\symbfit{r'}
    \label{eq:define}
\end{equation}

其中$i_k=1,2,3$。可以看出电多极矩的分量与电荷的角向分布密切相关，非零的电多极矩表明电荷体系是非球对称的。
\textcolor{blue}{那么反过来，任意阶电多极矩的分量均为零一定能推导出电荷分布是球对称的吗？}

\bigskip

\textcolor{red}{答案是否定的}，下面我们将构造出反例。

\bigskip

对于电势函数做多级展开，得到：

\begin{align}
    \phi(\symbfit{r})&=\frac{1}{4\pi\epsilon_0}\int\frac{\rho(\symbfit{r'})}{\left|\symbfit{r}-\symbfit{r'}\right|}\dd^3\symbfit{r'} \notag \\
    &=\frac{1}{4\pi\epsilon_0}\int\rho(\symbfit{r'})\sum_{l=0}^{\infty}\sum_{m=-l}^{l}\frac{4\pi}{2l+1}\frac{r'^l}{r^{l+1}}Y_{lm}(\theta,\phi)Y_{lm}^*(\theta',\phi')\dd^3\symbfit{r'} \notag\\
    &=\frac{1}{4\pi\epsilon_0}\frac{4\pi}{2l+1}\sum_{l=0}^{\infty}\sum_{m=-l}^{l}\left(\int\rho(\symbfit{r'})r'^lY_{lm}^*(\theta',\phi')\dd^3\symbfit{r'}\right)\frac{Y_{lm}(\theta,\phi)}{r^{l+1}}\\
    \label{eq:Multipole}
\end{align}

定义球谐函数表示的分量：

\begin{equation}
    p_{lm}=\int\rho(\symbfit{r'})r'^lY_{lm}^*(\theta',\phi')\dd^3\symbfit{r'}
\end{equation}

考虑到方程\eqref{eq:define}中分量$Q_{i_1i_2\cdots i_n}$只有$2n+1$个是独立的，且可以表示成$p_{lm},m=-l,-l+1,\cdots,l$
的线性组合，所以原问题相当于找一个电荷密度函数使得$\forall l,m\,\,p_{lm}=0$。考虑到球谐函数$Y_{lm}$的正交、完备、归一性，将
电荷密度$\rho(\symbfit{r})$以球谐函数作为基底展开得到：

\begin{equation}
    \rho(\symbfit{r})=\sum_{l=0}^{\infty}\sum_{m=-l}^{l}R_{lm}(r)Y_{lm}(\theta,\phi)
\end{equation}

那么原问题描述的情形对应于积分：

\begin{equation}
    p_{lm}=\int\left(\sum_{l'=0}^{\infty}\sum_{m'=-l}^{l'}R_{l'm'}(r)Y_{l'm'}(\theta,\phi)\right)r^lY_{lm}^*(\theta,\phi)\dd^3\symbfit{r}=0
\end{equation}

如果假设电荷分布在半径为$R$的球体内部，那么即有：

\begin{align}
    p_{lm}&=\sum_{l'=0}^{\infty}\sum_{m'=-l}^{l'}\left(\int_{0}^{R} R_{l'm'}(\symbfit{r})r^{l+2}\dd r\int Y_{l'm'}Y_{lm}^*\dd\Omega\right) \notag \\
    &=\int_{0}^{R} R_{lm}(\symbfit{r})r^{l+2}\dd r=0
    \label{eq:final}
\end{align}

方程\eqref{eq:final}是很容易满足的，例如假设电荷密度按照角向展开后只有$Y_{21}$分量，我们只需要构造$R(r)$满足：

\begin{equation}
    \int_{0}^{R} R_{lm}(\symbfit{r})r^{l+2}\dd r=0
\end{equation}

可以取：$R_{21}(r)=\frac{1}{r^4}\dfrac{\dd}{\dd r}r^6(R^6-r^6)=6R^6r-12r^{7}$，
此时有$\rho(\symbfit{r})=A(6R^6r-12r^{7})Y_{21}(\theta,\phi)$。
此电荷分布的任意阶极矩均为零。

\bigskip

以$\rho(\symbfit{r})=A(6R^6r-12r^{7})Y_{21}(\theta,\phi)$为例的一系列电荷构造
的电多极矩始终为零，说明在球体外部电场始终为零，而在球体内部对电荷体系做多级展开（内展开）
可知，球体内部电场非零而且电场具有丰富的角向结构。也就是说“电场线被囚禁在了有限区域内部”。

可以证明，对于其他形状的区域结论相同。这个问题本质上是判断在空间有限区域$\Omega$上，Poisson方程$\nabla^2\phi=\frac{\rho}{\epsilon}$
在边界条件$\phi(\partial\Omega)=0$是否有非球对称的解。从这个角度来说，答案就是显然的了，只要构造$\phi(\symbfit{r})$使得
电势在边界处为零且法向偏导数为零，且函数式显含角量，再代入Poisson方程即可得到电荷分布。（在区域$\Omega$之外，由于$r\to \infty,\phi \to 0$，根据静电场唯一性定理
可知，$\phi=0,\symbfit{E}=0$。）

\end{document}